\section{Introduction}
\label{sec:intro}

Training a large language model requires deciding, before the run starts, how to split a fixed compute budget $C$ between the number of model parameters $N$ and the number of training tokens $D$. Kaplan et al.~\citep{kaplan2020} first fit power laws for pre-training loss as a function of $N$ and $D$; Hoffmann et al.~\citep{hoffmann2022} revisited the question by training over 400 models and concluded that $N$ and $D$ should be scaled roughly in equal proportion, which became the widely cited rule of about 20 training tokens per parameter.

Several follow-up studies have questioned details of this result. Besiroglu et al.~\citep{besiroglu2024} attempted to replicate one of Hoffmann et al.'s three estimation approaches and reported inconsistencies in the published fit. Porian et al.~\citep{porian2024} and Pearce and Song~\citep{pearce2024} examined why the Kaplan and Hoffmann prescriptions differ, with the choice of how to count parameters (for example, whether embedding parameters are included) playing a role in the explanation. Most recently, Schaeffer et al.~\citep{schaeffer2025} refit the parametric loss on Chinchilla's data under three ways of computing the model sizes (the published counts, a formula assuming tied embeddings, and a best-fit variant of the formula; these counts differ by up to 15.2\%) and under four structured perturbations of the counts (multiplicative, additive, size-dependent bias, and log-normal noise). They report that the choice among the three interpretations does not meaningfully change the estimates or the roughly 20-to-1 token-to-parameter ratio, while additive or systematic errors can bias the fitted exponents and change the trend of the optimal ratio with compute.

These analyses rely on large training runs or on re-analysis of published data. It is less clear how such a robustness analysis behaves in the small-scale regime that a practitioner with limited compute can actually reproduce, where embedding and output-head parameters make up a large fraction of the model. We therefore ask:
\begin{itemize}
  \item \textbf{RQ1 (reproduction).} Do the sensitivities Schaeffer et al.\ report, in particular to additive offsets of the parameter count (counting or excluding embedding parameters acts like one) and to size-dependent bias, appear at small scale, where embedding parameters are a large share of the model?
  \item \textbf{RQ2 (boundary conditions).} As the fitting grid shrinks (fewer model sizes, fewer budgets), when do the conclusions stop being stable?
  \item \textbf{RQ3 (estimation approaches).} How do Hoffmann et al.'s three estimation approaches compare under each convention?
\end{itemize}

\paragraph{Contributions.} (i) A small-scale, fully open analogue of the parameter-counting analysis of Schaeffer et al., applying their four structured perturbations to our own fitted allocations with all three Hoffmann et al.\ estimation approaches. (ii) A grid-shrinking sweep that probes how the conclusions degrade with less data. (iii) Code and a recovery test on synthetic data with a known scaling law. \TODO{Edit contributions to match what the results actually support.}
